\chapter{TINJAUAN PUSTAKA}
\label{bab:tinjauan-pustaka}

\section{Regulasi Perpajakan Indonesia dan Keterkaitan Antar Pasal}

\section{Retrieval-Augmented Generation (RAG)}

\section{Temu-Kembali Hibrida (Lexical + Dense) dan Reciprocal Rank Fusion}

\section{Graph Retrieval-Augmented Generation (GraphRAG)}

\section{Basis Data Graf dan Neo4j}

\section{Embedding dan Model Bahasa Besar}

\section{Metrik Evaluasi Temu-Kembali}
% recall@k, precision@k; catatan RAGAS.

\section{Penelitian Terdahulu}
% CATATAN: hanya sitasi sumber yang benar-benar dibaca. Jangan mengarang referensi.
